\textbf{This is all stuff from the previous version, way to extensive here, all detailed analyes should be moved to a separate mapping paper:}

Results for gaussian rbf:
\begin{itemize}
  \item separate better than integrated polynomial (expected)
  \item n15 better than n10 better than n5 better than n3
  \item n15 and n10 too memory intensive (plot needed)
  \item onto coarser meshes is better than onto finer meshes (difference visible only mapping for very fine in meshes)
  \item separate between 1st and 2nd order
  \item integrated between 1st and 0th order
\end{itemize}

Results for thin plate splines rbf:
\begin{itemize}
  \item integrated and separate perform very similar 
  \item separate performs consistently slightly better than integrated.
\end{itemize}

Plots we want to do:
\begin{itemize}
  \item about the best way to go (complete picture): for one target mesh compare best local-rbf (gaussian), global-rbf (tps), nn, np (For each geometry: 12 series)
    Geometry: Turbine only (Aorta is too fine)
    from: 0.01 (1.5k V) - 0.0075 (2.7k) -  (0.004 TODO ) - 0.0025 (25k, no tps)
    to: 0.005 (6.2k)
    Goal: show errors of best mappings including tps (user perspective) (small meshes only to show how tps performs for cases where it is available)

  \item all geometries, pick one target mesh each: nn, np, best rbf
    Geometry: ALL in on plot
    from: 0.01 and finer
    to: 2nd finest
    Goal: show errors of best mappings for all geometries excluding tps (user perspective)

  \item pick one geometry, one target mesh: nn, np, gaussian-nX separate
    Geometry: Turbine
    from: 0.01 and finer
    to: 2nd finest
    Y-axis: plot time of computeMapping and map and peak memory
    Goal: show peak memory usage and computeMapping and map time of best mappings (user perspective)

  \item pick one geometry and target mesh: compare gaussian-nX on vs separate for varying X (8 series)
    Geometry: Turbine
    from: 0.01 and finer
    to: 2nd finest
    Goal: show errors (and hopefully runtime) of best mappings (user perspective)
    % Goal: Options for local-rbf you should not choose and why

  \item pick one geometry, fixed src and tgt meshes, varying rank counts: nn, np, gaussian-nX separate (best error)
    Y-axis: plot time of computeMapping max + map max (upper limit) 
    Goal: show strong scalability of best mappings (user perspective)
    Participant ranks on A are not so important. Part B is the one that computes the mapping.
    Use micro partition (to start with) 1 Node for Part A, 1,2,4,8,16 Nodes for Part B. Scale with whole nodes.
    % Strong scaling is defined as how the solution time varies with the number of processors for a fixed total problem size. 

  \item pick one geometry, one target mesh, varying rank counts: nn, np, gaussian-nX separate (best error)
    Y-axis: plot peak memory 
    Goal: show the memory overhead for using more ranks

  % Decide after the above. Was shown in the past with the uniform plate. No real added value here.
  \item varying rank count, varying meshes: nn, np, gaussian-nX separate
    Geometry: uniform plate with uniform partitioning
    Choose from and to mesh for serial case. Then double the vertex count of both meshes an double the rank count. Keep problem size per rank constant.
    Y-axis: plot time of computeMapping max + map max (upper limit) 
    Goal: show weak scalability of best mappings (user perspective)
    % Weak scaling is defined as how the solution time varies with the number of processors for a fixed problem size per processor.
\end{itemize}

\todo{FS - evaluate\_error}
\todo{Kyle - fix aste connectivity and partitioning}
\todo{Kyle - meshes}
\todo{FS - (PETSc on) SuperMUC}
\todo{FS - create configs and jobscripts}
\todo{FS - run first accuracy tests (NN, NP)}
% RBF postponed

\paragraph{Numerical Accuracy}

have a look at Bernhard's thesis, Section 4.2.4
also have a look at Linder et al. 2017 "Radial Basis Function Interpolation for Black-Box Multi-Physics Simulations"
also Florian's thesis

\begin{itemize}
  \item which preCICE version to test: v2.0.2
  \item which ASTE version? - create a tag

  \item how to configure RBF, compare different settings (basis function, support radius, shape parameter)\\
    basis functions: gaussian, thin-plate-splines, tps-compact-c2, multi-quadratic\\
    support radius?? trial and error ( also look at extras/rbfShape )

  \item classical convergence plots
    challenge: we can refine the input and the output mesh

  \item spatial error plot, screenshot of ParaView, only for one specific (interesting) setting

  \item which functions? some $sin(x+y+z)$?, also some try and error, we want to avoid edge cases
    Lindner Thesis - Rosenbrock, eggholder, step function
\end{itemize}

\paragraph{Runtime and Scalability}

\begin{itemize}
  \item target system: SuperMUC-NG.
  \item ideally scalability until 10,000 ranks per side.
  \item for approximate numbers of vertices, you could check BU's thesis, NN and NP (not RBF, has changed much since then), ASTE tests and "Ateles Cube"
  \item only measure mapping parts ("computeMapping", "map"), optional: sub-stages (but not so interesting for a user)
\end{itemize}
